\section{Deep-Dream}
\label{sec:deep-dream}

\subsection{记忆表示}

设 $\mathcal D_t$ 是截至处理时刻 $t$ 已处理的文档集合。对 $D\in\mathcal D_t$，设 $\mathcal E(D)$ 是从该文档不可修改的原文片段中切分出的 Episode。对 $e_i\in\mathcal E(D_i)$，定义一条存储观测为
\[
o_i=(D_i,e_i,\tau_i,x_i,\ell_i),
\]
其中 $\tau_i$ 是处理时间，$x_i$ 是抽取出的内容和属性，$\ell_i$ 是它在 $D_i$ 中的精确原文位置。写入 $t$ 条观测后的观测库为 $O_t=\{o_1,\ldots,o_t\}$。持久化状态为
\[
M_t=(F_t,O_t,G_t),
\]
其中 $F_t$ 保存概念实体族和候选身份，$G_t$ 保存关系、邻域与版本指针。每条观测都携带文档和 Episode 标识 $(D_i,e_i)$，用它们可以重新定位原文。

一次写入状态转移表示为
\begin{equation}
\begin{aligned}
F_{t+1}&=\mathsf{Align}(F_t,x_{t+1}),\\
O_{t+1}&=O_t\cup\{o_{t+1}\},\\
G_{t+1}&=\mathsf{Link}(G_t,F_{t+1},o_{t+1}),\\
M_{t+1}&=(F_{t+1},O_{t+1},G_{t+1}).
\end{aligned}
\label{eq:state-transition}
\end{equation}
这里 $\mathsf{Align}$ 要么把 $x_{t+1}$ 分配到已有实体族，要么先保留一个新的候选实体族；$\mathsf{Link}$ 在实体族、观测和版本之间增加或重定向指针。两种操作都不会修改 $\ell_{t+1}$ 所指向的原文。由于版本指针保留了早期状态，沿版本链定义的投影 $\pi_t$ 可以从 $M_{t+1}$ 恢复 $M_t$；所以状态更新是增加链接，而不是删除早期观测。

这种表示把“找到”与“读取”分开。解析实体族或关系属于索引查询，打开对应的源文本则是单独的读取操作。关系边帮助智能体找到相连的观测，但支撑这条边的陈述仍保存在原文中。处理时间 $\tau_i$ 是图谱唯一的状态轴。原文中的事件日期可以放在 $x_i$ 中并参与回答，但不会产生第二套版本时钟。

\subsection{写入记忆}

写入遵循\cref{eq:state-transition}中的状态转移：按照原文位置切分文档中的 Episode；从每个 Episode 中抽取实体和关系；将候选身份与已有概念实体族比较；最后把确认的链接和支撑它们的观测一起写入。每条链接都能沿路径回到具体文档片段。

首次提及无法确定实体身份时，Deep-Dream 会先保留多个候选实体族。后续观测和关系邻域提供足够证据后，再合并兼容的实体族或重定向指针。这个操作只修改 $F_t$ 和 $G_t$，不修改 $O_t$ 中的观测。因此，暂时拆分可以在之后修正；相反，过早合并可能把两段历史混在同一个检索入口下。

\subsection{即时对齐与延迟实体对齐}

实验比较两种写入时机，它们使用相同的状态表示。 \textbf{即时对齐（Base）} 在每个窗口到达时执行 $\mathsf{Align}$；\textbf{延迟对齐（Align+）} 先追加 $o_i$ 和候选链接，等周围窗口处理完后再解析兼容的实体族。两者的区别在于身份决策何时发生，而不是图谱结构或原文记录不同。

延迟对齐在做合并之前使用更多上下文。它可以把兼容候选放在一起处理，减少重复工作，同时保留支撑每次决策的源文本。第~\cref{sec:temporal-benchmark} 节的配对结果量化了这种取舍。

\subsection{读取记忆}
\label{sec:answer-paths}

设 $\mathcal S_t$ 是 $M_t$ 中全部不可修改的源文本片段。对问题 $q$，智能体从 $C_0(q)=\varnothing$ 开始，选择动作
\[
a_{k+1}=\rho\bigl(q,M_t,C_k(q),u_k\bigr)
\in\bigl(\mathcal S_t\setminus C_k(q)\bigr)\cup\{\mathsf{stop}\},
\]
其中 $u_k$ 是智能体当前的不确定性。如果 $a_{k+1}$ 是一个源文本片段，就打开它并令
\[
C_{k+1}(q)=C_k(q)\cup\{a_{k+1}\};
\]
如果 $a_{k+1}=\mathsf{stop}$，读取过程结束。用归纳法——每一步只把 $\mathcal S_t$ 中的一个片段加入 $C_k$——可得 $C_k(q)\subseteq\mathcal S_t$ 且 $C_k(q)\subseteq C_{k+1}(q)$。这个公式只记录打开源文本的动作；词法、语义、概念和关系查询会更新 $u_k$，但不会把片段加入 $C_k$。这就是渐进读取的形式化定义；何时停止由策略决定。

提问时，智能体从当前最有用且成本较低的线索——词法、语义、概念或关系——开始，打开对应 Episode 或源文本，在证据不完整时继续扩展：只问一个事实时，可能一段原文就够了；涉及更新或多个文档时，则可能需要沿关系和版本链接继续读取。

我们比较三种读取方式：
\begin{enumerate}
    \item \textbf{固定检索。} 只运行一次词法/语义/图检索，把一份排序后的记忆记录和检索片段交给答题模型；列表不能再修改，也不会与源文件核对。
    \item \textbf{工具辅助读取。} 把只读的搜索、概念、关系和片段工具交给答题智能体；它可以多次调用，并用一次结果组织下一次查询，但答题模型拿到的是工具返回的记录。
    \item \textbf{原文读取。} 用索引选择文档组，再打开原始文件、读取精确片段，并且只提交通过原文门禁的片段。
\end{enumerate}

源文本检查器确认引用的标识符确实由工具返回，并且仍然指向原文。它保证引用来源可追溯，但仍需要智能体判断该片段是否真正支持答案。

\subsection{为什么这种表示能够保留有用细节}
\label{sec:information-view}

设 $Z$ 是观测对应的潜在实体身份。由于 $\pi_t(M_{t+1})=M_t$，较新的状态包含早期状态，以及新观测和修订后的链接。因此
\begin{equation}
\begin{aligned}
H(Z\mid M_t)-H(Z\mid M_{t+1})
&=I(Z;M_{t+1}\mid M_t)\\
&\ge 0.
\end{aligned}
\label{eq:state-uncertainty}
\end{equation}
第一个等式成立是因为 $M_t$ 可以由 $M_{t+1}$ 恢复——这是由保留的版本链接提供的表示不变量，并不是任意图边都能推出的。如果新观测不含有区分 $Z$ 的信息，差值就是零；如果它能区分候选实体族，条件熵就会下降。这是保留状态的结构性保证，不是已经测量得到的熵增益。

对于问题 $q$，设 $A$ 是用于定位候选的紧凑记录，$E_C$ 是从候选集合 $C$ 中打开的源文本，$Y$ 是与答案有关的事实。链式法则给出
\begin{equation}
I(Y;A,E_C\mid Q=q)=I(Y;A\mid Q=q)+I(Y;E_C\mid A,Q=q).
\label{eq:two-rate-code}
\end{equation}
这就是链式法则。第一项表示定位候选时使用的信息，第二项表示读取源文本后获得的额外信息。Deep-Dream 压缩的是检索描述，而不是事实的唯一副本。

设 $R_q$ 表示相关源文本位于已打开集合 $C_K(q)$ 中，设 $V_q$ 表示最终答案正确，即 $V_q=\{\widehat Y=Y\}$。由全概率公式，
\begin{equation}
\begin{aligned}
\Pr[\widehat Y=Y\mid q]
&=\Pr[R_q\mid q]\Pr[V_q\mid R_q,q]\\
&\quad+\Pr[\neg R_q\mid q]\Pr[V_q\mid\neg R_q,q]\\
&\ge \Pr[R_q\mid q]\Pr[V_q\mid R_q,q].
\end{aligned}
\label{eq:answer-decomposition}
\end{equation}
概念和关系链接主要影响第一项，原文读取和引用检查影响第二项。这个不等式不是性能估计，只是把正确回答分解为“找到相关证据”和“读取后正确作答”两个环节。
